\chapter{Исследование алгоритмов управления УСП} \label{chapt4}

\section{Структурная схема алгоритма управления УСП} \label{sect4_1}

На рисунке \ref{img:SchemeControl} представлены варианты схемы системы управления СП.

\begin{figure}[ht]
    \center{
        \hfill
        \subcaptionbox[List-of-Figures entry]{\label{img:SchemeControlTaskSpace}} {\includegraphics[width=\linewidth]{SchemeControlTaskSpace}}%
        \hfill
        \subcaptionbox{\label{img:SchemeControlSeparate}} {\includegraphics[width=\linewidth]{SchemeControl}}
        \hfill
    }

    Варианты схемы системы управления СП
    \caption{Схема системы управления СП} % Этот текст попадает в названия рисунков в списке рисунков
    \label{img:SchemeControl}
\end{figure}

Условно обозначим вариант системы управления, изображенный на рисунке \ref{img:SchemeControlSeparate} как <<система 1>>, а вариант, изображенный на рисунке \ref{img:SchemeControlTaskSpace} – <<система 2>>:
	С состав системы 1 входят:
\begin{itemize}
	 \item блок ограничительной характеристики (ОХ), осуществляющий проверку корректности значений задаваемых координат платформы с целью обеспечения возможности их достижения с помощью данной системы подвижности;
	\item блок решения обратной задачи кинематики (ОЗК), осуществляющий преобразование прошедших предварительную обработку координат  платформы в углы поворота КШМ $q_1^*, q_2^*, q_3^*$
	\item локальные системы управления (ЛСУ) углами поворота КШМ, включающие: ПИД-регуляторы, блоки ограничения управляющих воздействий (БОУВ), преобразователи частоты (ПЧ), датчики угла (ДУ) поворота мотор-редукторов, мотор-редукторы (МР) приводных звеньев;
	\item объект управления (ОУ).
\end{itemize}
	В состав системы 2 дополнительно входят:
\begin{itemize}
	\item блок прямой задачи кинематики (ПЗК), осуществляющий преобразование измеряемых обобщенных координат $q_1, q_2, q_3$ в текущие координаты платформы
	\item регулятор положения платформы (РПП), производящий расчет скоростей  ,  (частот вращения $f1, f2, f3$) МР исходя из ошибок позиционирования  $\Delta\alpha, \Delta\beta, \Delta z_p$.
БОУВ как для системы 1, так и для системы 2 дополнительно осуществляют масштабирование выходов регуляторов $u_1, u_2, u_3$ в частоты $f_1, f_2, f_3$ для корректности задания частоты вращения, передаваемой на  ПЧ.
\end{itemize}
	Перейдем к описанию базовых алгоритмов представленной на рисунке \ref{img:SchemeControl} системы управления.

\section{Блок ограничительной характеристики} \label{sect4_2}

\begin{figure}
	\center
	\caption{К алгоритму <<мягкой>> срезки задаваемых углов платформы}
	\includegraphics[width=0.75\linewidth]{zoneRestriction}
	\label{img:zoneRestriction}
\end{figure}

В следящем режиме работы заданные координаты платформы необходимо ограничивать перед передачей на блок решения ОЗК, т.е. значения заданных координат X, находящиеся за границами возможностей платформы, должны быть <<притянуты>> к границам ограничительной характеристики. Для определения точки, ближайшей к границе ограничительной зоны, используется алгоритм, проиллюстрированный на рисунке \ref{img:zoneRestriction}.
	Алгоритм функционирует следующим образом. Строится ограничительная характеристика при $z_p = const$. Если заданная точка $X2$ находится за пределами области допустимых перемещений, производится поиск допустимого решения на отрезке $OX2$ путем его деления в пропорции золотого сечения. На каждом шаге данного алгоритма решается ОЗК для очередной точки внутри малого отрезка X2’X2’’ после чего проверяется условие останова. При достижении заданной точности искомая точка $X2*$ передается на блок ОЗК.

\section{Регулятор положения платформы} \label{sect4_3}

Для реализации контуров положения ЛСУ системы 1, замкнутых обратными связями, целесообразно применить ПД-регулятор с ограничением по управляющим воздействиям (скоростям) вследствие простоты его реализации и приемлемого качества переходных процессов. Путем настройки коэффициентов регулятора \cite{Plohotnikov, Voevoda} достигаются требуемые характеристики переходных процессов в контурах ЛСУ при реакции на типовые сигналы задания \cite{ChLewin} (как для каждого контура в отдельности, так и при совместной работе в пределах ограничительной характеристики).

Основная формула ПД-регулятора с независимыми каналами выглядит следующим образом:
\begin{equation}
\label{eq:ControllerSystem1}
\begin{aligned}
	\begin{pmatrix}
	u_1 \\ u_2 \\ u_3
	\end{pmatrix} =
	\begin{pmatrix}
	K_{p1} & 0 & 0 \\
	0 & K_{p2} & 0 \\
	0 & 0 & K_{p3} \\
	\end{pmatrix}
	\begin{pmatrix}
	e_{q1} \\ e_{q2} \\ e_{q3}
	\end{pmatrix} +
	\begin{pmatrix}
	K_{d1} & 0 & 0 \\
	0 & K_{d2} & 0 \\
	0 & 0 & K_{d3} \\
	\end{pmatrix}
	\begin{pmatrix}
	\dot{e}_{q1} \\ \dot{e}_{q2} \\ \dot{e}_{q3}
	\end{pmatrix}
\end{aligned}
\end{equation}

Для системы 2 многомерный ПД-регулятор положения платформы \cite{Hernandez2009} строится в соответствии с учетом влияния каждой из компонент вектора ошибки на вектор $\mathbf{u}$:

\begin{equation}
\label{eq:ControllerSystem2}
\begin{aligned}
	\begin{pmatrix}
	u_1 \\ u_2 \\ u_3
	\end{pmatrix} =
	\begin{pmatrix}
	K_{\alpha}^p & K_{\beta1}^p & K_h^p \\
	-K_{\alpha}^p & K_{\beta1}^p & K_h^p \\
	0 & -K_{\beta2}^p & K_h^p \\
	\end{pmatrix}
	\begin{pmatrix}
	e_{\alpha} \\ e_{\beta} \\ e_{Zp}
	\end{pmatrix} +
	\begin{pmatrix}
	K_{\alpha}^d & K_{\beta1}^d & K_h^d \\
	-K_{\alpha}^d & K_{\beta1}^d & K_h^d \\
	0 & -K_{\beta2}^d & K_h^d \\
	\end{pmatrix}
	\begin{pmatrix}
	\dot{e}_{\alpha} \\ \dot{e}_{\beta} \\ \dot{e}_{Zp}
	\end{pmatrix}
\end{aligned}
\end{equation}

Для оценки качества переходных и установившихся процессов в системах 1 и 2 воспользуемся критерием минимума среднего квадрата ошибки:

\begin{equation}
\label{eq:KriteriaControlQuality}
\begin{aligned}
	J = \frac{\displaystyle \int_0^{T_{\text{уст}}} \left(c_1e_{\alpha}^2 + c_2e_{\beta}^2 + c_3e_{z_p}^2\right)dt}{T_{\text{уст}}} \rightarrow \min
\end{aligned}
\end{equation}

	где $c_1, c_2, c_3$ – нормирующие коэффициенты, $e_{\alpha}, e_{\beta}, e_{Zp}$ – ошибки регулирования по $\alpha, \alpha, Zp$ соответственно.
Данный критерий выбран исходя из основного назначения СП в составе тренажерного комплекса -– движение по заданной траектории, формируемой математической моделью техники. Математическая модель АТ и СТ отвечает за корректный расчет параметров акселерационных воздействий на рабочем месте оператора, на основании которых формируется траектория движения СП. К системе управления СП при этом предъявляются требования, описанные в \ref{sect1_4}.

\begin{table} [htbp]%
    \centering
	\caption{Значение критерия для синусоидальных воздействий различной частоты}%
	\label{tbl:KriteriaFreqTable}
    \renewcommand{\arraystretch}{1.5}
    \begin{SingleSpace}
	\begin{tabulary}{\textwidth}{@{}>{\zz}L >{\zz}L >{\zz}L >{\zz}L >{\zz}L@{}}
        \toprule     %%% верхняя линейка
    	Частота $f$, Гц	& Критерий $J_1$ & Критерий $J_2$\\
        \midrule
    	0,1 &	0,0083 &	0,0096 \\
		0,2 &	0,0096 &	0,0109 \\
		0,3 &	0,0112 &	0,0126 \\
		0,4 &	0,0130 &	0,0143 \\
		0,5 &	0,0150 &	0,0159 \\
		0,6 &	0,0171 &	0,0174 \\
		0,7 &	0,0192 &	0,0188 \\
		0,8 &	0,0217 &	0,0205 \\
		0,9 &	0,0260 &	0,0251 \\
        \bottomrule
	\end{tabulary}%
   	\end{SingleSpace}
\end{table}

Рассмотренные выше варианты решения задачи управления положением платформы были испытаны на стенде тренажера БТР-82А.

В таблице 1 приведены результаты испытаний по отработке простых синусоидальных воздействий следующего вида:

Для амплитуды колебаний  $\alpha = \beta = 10^o$, значение критерия $J$, удовлетворяющее перечисленным выше требованиям к системе, находится в пределах $J < 0,03$.

Для сравнения систем 1 и 2 по критерию $J$ были подобраны эквивалентные настройки регуляторов.

\begin{figure}
	\center
	\caption{Сравнение отработки задающих воздействий систем 1 и 2}
	\label{img:testKriteriaControlSystem}
	\includegraphics[width=\linewidth]{testKriteriaControlSystem}
\end{figure}

Сравнивая кривые для системы 1 и 2 и данные таблицы \ref{tbl:KriteriaFreqTable} можно прийти к выводу, что отработка задающих воздействий в следящем режиме в обоих случаях происходит схожим образом, однако на больших частотах сигнала задания система 2 лучше компенсирует отклонения ошибки по скорости за счет влияния коэффициентов обратной связи. Система 2 также выгодно отличается большими возможностями настройки. Так, можно уменьшить коэффициенты по осям $\alpha-\beta$ для более плавной отработки низкочастотных сигналов или обнулить коэффициенты при $h$ для задействования большего объема рабочей зоны.

\section{Методика настройки параметров алгоритма управления УСП} \label{sect4_4}

Для первичной настройки коэффициентов регуляторов в \ref{eq:ControllerSystem1} и \ref{eq:ControllerSystem2} целесообразно использовать динамическую модель системы подвижности.

Структурная схема динамической модели, проиллюстрированная на рисунке \ref{img:FullDynamicalModel} аналогична уже рассмотренной выше схеме \ref{img:SchemeControl}

В данной модели вводятся следующие допущения:
\begin{itemize}
\item платформа считается абсолютно жесткой, упругие связи между звеньями механизма считаются несущественными;
\item влияние масс и моментов инерции звеньев механизма (кривошип, шатун, стойка) не рассматривается;
\item люфт передач, транспортные задержки считаются несущественными;
\item частота дискретизации цифрового ПИД-регулятора считается постоянной.
\end{itemize}

\begin{figure}[h!]
\caption{Полная динамическая модель СП}
\label{img:FullDynamicalModel}
\center
\includegraphics[width=\linewidth]{control_system_complicated}
\end{figure}

В дополнение к описанным выше блокам в данной модели присутствуют следующие элементы:
\begin{itemize}
\item модель системы управления приводным звеном (СУ ПЗ), построенная на базе передаточной функции двигателя постоянного тока;
\item модуль решения прямой задачи динамики (ПЗД). решение ПЗД описано в главе \ref{chapt3}.
\end{itemize}

Модель СУ ПЗ показана на рисунке \ref{img:ServoSystemModel}. Структура данной системы известна \cite{Najjari2012, Firago, Ignatova, 2004a, KolganovControlTheory} и включает двигатель постоянного тока, ПИ-регулятор тока/момента, ПИ-регулятор скорости, соединенные каскадно и реализующие подчиненное управление электроприводом.

\begin{figure}[!htbp]
\center
\caption{Модель приводного звена на базе сервопривода}
\label{img:ServoSystemModel}
\includegraphics[width=0.9\linewidth]{DCmotorWithController}
\end{figure}

Процесс настройки коэффициентов регулятора в целом достаточно хорошо изучен \cite{Wang2001,Yang}, и его алгоритм проиллюстрирован на рисунке \ref{img:ControlSystemTuning}:

\begin{figure}[h!]
\center
\caption{Алгоритм настройки параметров КД}
\includegraphics[width=0.7\linewidth]{controlSystemTuning}
\label{img:ControlSystemTuning}
\end{figure}
